\documentclass[12pt]{article}
\usepackage[utf8]{inputenc}
\usepackage[a4paper, margin=2cm]{geometry}
\usepackage{amsthm}
\usepackage{amsmath, amssymb}
\usepackage{newtxtext, newtxmath}
\usepackage{tcolorbox}
\usepackage{titlesec}
\usepackage{xcolor}
\definecolor{lightblue}{RGB}{135, 206, 250}
\tcbuselibrary{skins,theorems}

\titleformat{\section}[block]{\Large\bfseries\centering}{}{0pt}{}
\titleformat{\subsection}[block]{\bfseries}{\thesubsection}{2pt}{}

\newtcbtheorem[no counter]{proposition}{\large\hspace{-8pt}Proposition}{
  enhanced,
  colback=white,
  colframe=lightblue,
  left=12pt,
  right=12pt,
  top=24pt,
  bottom=24pt,
  fonttitle=\bfseries
}{prop}
\newtcbtheorem[]{question}{\large\hspace{-8pt}Question :}{
  enhanced,
  colback=white,
  colframe=black,
  left=12pt,
  right=12pt,
  top=24pt,
  bottom=24pt,
  fonttitle=\bfseries
}{qst}
\newenvironment{Answer}{
  \vspace{12pt}
  {\large\textbf{\textit{\hspace{-12pt}Answer:}}}\par
  \vspace{8pt}
} {\vspace{24pt}}
\newenvironment{Proof}{
  \vspace{12pt}
  {\large\textbf{\textit{\hspace{-12pt}Proof:}}}\par
  \vspace{8pt}
} {\vspace{24pt}}


\title{DCHw1}
\author{Miki-Riako }
\date{\today}
\begin{document}

\begin{question}{}{}
  1-1 Binary to decimal conversion.

  (1) 1100110; (2) 11011011
\end{question}
\begin{Answer}
\noindent
\begin{minipage}{0.5\textwidth}
\begin{flalign*}
(1)\;\;1100110_{(2)} &= 2^6+2^5+2^2+2^1 &&\\
  &= 64+32+4+2 &&\\
  &= 102
\end{flalign*}
\end{minipage}\begin{minipage}{0.5\textwidth}
\begin{flalign*}
(2)\;\;11011011_{(2)} &= 2^7+2^6+2^4+2^3+2^1+2^0 &&\\
  &= 128+64+16+8+2+1 &&\\
  &= 219
\end{flalign*}
\end{minipage}
\end{Answer}

\begin{question}{}{}
  1-2 Decimal to binary conversion.

  (1) 31; (2) 64 
\end{question}
\begin{Answer}
\noindent
\begin{minipage}{0.5\textwidth}
\begin{flalign*}
(1)\;\;31 &= 32-1 &&\\
  &= 2^5-1 &&\\
  &= 11111_{(2)}
\end{flalign*}
\end{minipage}\begin{minipage}{0.5\textwidth}
\begin{flalign*}
(2)\;\;64 &= 2^6 &&\\
  &= 1000000_{(2)} &&\\
  \;
\end{flalign*}
\end{minipage}
\end{Answer}

\begin{question}{}{}
  1-3 Decimal to hexadecimal conversion.

  (1) 96; (2) 789
\end{question}
\begin{Answer}
\noindent
\begin{minipage}{0.5\textwidth}
\begin{flalign*}
(1)\;\;96 &= 6\times 16+0 &&\\
  &= 60_{(16)}
\end{flalign*}
\end{minipage}\begin{minipage}{0.5\textwidth}
\begin{flalign*}
(2)\;\;789 &= 3\times 16^2+1\times 16+5&&\\
  &= 315_{(16)}
\end{flalign*}
\end{minipage}
\end{Answer}

\begin{question}{}{}
  1-4 Binary to hexadecimal conversion.

  (1) 10110110; (2) 101011
\end{question}
\begin{Answer}
\noindent
\begin{minipage}{0.5\textwidth}
$(1)\;\;10110110_{(2)} = B6_{(16)}$
\end{minipage}\begin{minipage}{0.5\textwidth}
$(2)\;\;101011(2) = 2B_{(16)}$
\end{minipage}
\end{Answer}

\begin{question}{}{}
  1-5 Hexadecimal to binary conversion.

  (1) F4; (2) A0B
\end{question}
\begin{Answer}
\noindent
\begin{minipage}{0.5\textwidth}
$(1)\;\;F4_{(16)} = 11110100_{(2)}$
\end{minipage}\begin{minipage}{0.5\textwidth}
$(2)\;\;A0B_{(16)} = 101000001000_{(2)}$
\end{minipage}
\end{Answer}

\begin{question}{}{}
  1-6 Decimal to BCD code conversion.

  (1) 6; (2) 47
\end{question}
\begin{Answer}
\noindent
\begin{minipage}{0.5\textwidth}
$(1)\;\;6 = 0110_{(BCD)}$
\end{minipage}\begin{minipage}{0.5\textwidth}
$(2)\;\;47 = 01000111_{(BCD)}$
\end{minipage}
\end{Answer}

\begin{question}{}{}
  1-7 Binary primitive, inverse, and complementary codes with the length of 5.

  (1) +1100; (2) -1100
\end{question}
\begin{Answer}
\begin{table}[h]
\centering
\begin{tabular}{|c|c|c|c|}
\hline
& Primitive & Inverse & Complementary \\ \hline
(1) +1100 & 01100 & 01100 & 01100 \\ \hline
(2) -1100 & 11100 & 10011 & 10100 \\ \hline
\end{tabular}
\end{table}
\end{Answer}

\begin{question}{}{}
  1-8 Primitive, inverse, and complementary codes of binary.

  (1) 01110; (2) 10101  
\end{question}
\begin{Answer}
\begin{table}[h]
\centering
\begin{tabular}{|c|c|c|c|}
\hline
& Primitive & Inverse & Complementary \\ \hline
(1) 01110 & 01110 & 01110 & 01110 \\ \hline
(2) 10101 & 10101 & 11010 & 11011 \\ \hline
\end{tabular}
\end{table}
\end{Answer}

\begin{question}{}{}
  1-9 Complementary codes with the length of 8 for decimal.

  (1) +15; (2) -15
\end{question}
\begin{Answer}
\begin{minipage}{0.5\textwidth}
(1) 00001111
\end{minipage}\begin{minipage}{0.5\textwidth}
(2) 11110001
\end{minipage}
\end{Answer}

\begin{question}{}{}
  
\end{question}


\end{document}